% GENERATED FILE. Edit canonical lessons, then run npm run generate:books.
% canonical-source-hash: fnv1a64:eeec99d739755dd3
% canonical-lessons: TA-C132-kocu, TA-C132-mukam, TA-C132-vay, TA-C132-nakku, TA-C132-viral

\chapter{Mosquito, Face, Mouth, Tongue, Finger}
\label{ch:ta-mosquito-face-mouth-tongue-finger}
\hlchaptermodality{132}

\begin{chapteropening}
\textbf{By the end of this chapter:} \emph{I can say a mosquito, the face, the mouth, the tongue and a finger.}

Complete the last lesson of chapter 132: i can say a mosquito, the face, the mouth, the tongue and a finger.
\end{chapteropening}

\section[kocu]{\ta{கொசு} (kocu) — a mosquito}

\label{lesson:TA-C132-kocu}



\begin{quote}
Before the new one: say the Tamil for a peacock, then the Tamil for an ant.
\end{quote}

\subsection*{You'll want to know: \ta{கொசு}}
\textbf{\ta{கொசு}} — \emph{kocu} — \textquotedblleft{}a mosquito\textquotedblright{}.

\begin{grammarlens}[title={What you've built}]
One more word for this chapter.
\end{grammarlens}

\subsection*{Your turn}
\begin{itemize}
\raggedright
  \item \emph{Say it:} \emph{kocu}
  \item \emph{Say it:} \emph{kocu}, once more
  \item \emph{Say it:} say \emph{kocu}
  \item \emph{Recall:} say the Tamil for a peacock, then the Tamil for an ant, then say \emph{kocu} again
\end{itemize}

\begin{tcolorbox}[breakable,colback=teal!4,colframe=teal!35!black,title={Before you move on}]
What does \ta{கொசு} mean? (\textquotedblleft{}A mosquito\textquotedblright{}.) Say it once more.
\end{tcolorbox}
\section[mukam]{\ta{முகம்} (mukam) — the face}

\label{lesson:TA-C132-mukam}



\begin{quote}
Before the new one: say the Tamil for an ant, then the Tamil for a mosquito.
\end{quote}

\subsection*{You'll want to know: \ta{முகம்}}
\textbf{\ta{முகம்}} — \emph{mukam} — \textquotedblleft{}the face\textquotedblright{}.

\begin{grammarlens}[title={What you've built}]
One more word for this chapter.
\end{grammarlens}

\subsection*{Your turn}
\begin{itemize}
\raggedright
  \item \emph{Say it:} \emph{mukam}
  \item \emph{Say it:} \emph{mukam}, once more
  \item \emph{Say it:} say \emph{mukam}
  \item \emph{Recall:} say the Tamil for an ant, then the Tamil for a mosquito, then say \emph{mukam} again
\end{itemize}

\begin{tcolorbox}[breakable,colback=teal!4,colframe=teal!35!black,title={Before you move on}]
What does \ta{முகம்} mean? (\textquotedblleft{}The face\textquotedblright{}.) Say it once more.
\end{tcolorbox}
\section[vāy]{\ta{வாய்} (vāy) — the mouth}

\label{lesson:TA-C132-vay}



\begin{quote}
Before the new one: say the Tamil for a mosquito, then the Tamil for the face.
\end{quote}

\subsection*{You'll want to know: \ta{வாய்}}
\textbf{\ta{வாய்}} — \emph{vāy} — \textquotedblleft{}the mouth\textquotedblright{}.

\begin{grammarlens}[title={What you've built}]
One more word for this chapter.
\end{grammarlens}

\subsection*{Your turn}
\begin{itemize}
\raggedright
  \item \emph{Say it:} \emph{vāy}
  \item \emph{Say it:} \emph{vāy}, once more
  \item \emph{Say it:} say \emph{vāy}
  \item \emph{Recall:} say the Tamil for a mosquito, then the Tamil for the face, then say \emph{vāy} again
\end{itemize}

\begin{tcolorbox}[breakable,colback=teal!4,colframe=teal!35!black,title={Before you move on}]
What does \ta{வாய்} mean? (\textquotedblleft{}The mouth\textquotedblright{}.) Say it once more.
\end{tcolorbox}
\section[nākku]{\ta{நாக்கு} (nākku) — the tongue}

\label{lesson:TA-C132-nakku}



\begin{quote}
Before the new one: say the Tamil for the face, then the Tamil for the mouth.
\end{quote}

\subsection*{You'll want to know: \ta{நாக்கு}}
\textbf{\ta{நாக்கு}} — \emph{nākku} — \textquotedblleft{}the tongue\textquotedblright{}.

\begin{grammarlens}[title={What you've built}]
One more word for this chapter.
\end{grammarlens}

\subsection*{Your turn}
\begin{itemize}
\raggedright
  \item \emph{Say it:} \emph{nākku}
  \item \emph{Say it:} \emph{nākku}, once more
  \item \emph{Say it:} say \emph{nākku}
  \item \emph{Recall:} say the Tamil for the face, then the Tamil for the mouth, then say \emph{nākku} again
\end{itemize}

\begin{tcolorbox}[breakable,colback=teal!4,colframe=teal!35!black,title={Before you move on}]
What does \ta{நாக்கு} mean? (\textquotedblleft{}The tongue\textquotedblright{}.) Say it once more.
\end{tcolorbox}
\section[viral]{\ta{விரல்} (viral) — a finger}

\label{lesson:TA-C132-viral}



\begin{quote}
Before the new one: say the Tamil for the mouth, then the Tamil for the tongue.
\end{quote}

\subsection*{You'll want to know: \ta{விரல்}}
\textbf{\ta{விரல்}} — \emph{viral} — \textquotedblleft{}a finger\textquotedblright{}.

\begin{grammarlens}[title={What you've built}]
One more word for this chapter.
\end{grammarlens}

\subsection*{Your turn}
\begin{itemize}
\raggedright
  \item \emph{Say it:} \emph{viral}
  \item \emph{Say it:} \emph{viral}, once more
  \item \emph{Say it:} say \emph{viral}
  \item \emph{Recall:} say the Tamil for the mouth, then the Tamil for the tongue, then say \emph{viral} again
\end{itemize}

\begin{tcolorbox}[breakable,colback=teal!4,colframe=teal!35!black,title={Before you move on}]
What does \ta{விரல்} mean? (\textquotedblleft{}A finger\textquotedblright{}.) Say it once more.
\end{tcolorbox}
